% !TEX root = ../main.tex

\begin{abstract}
Evolutionary feature-selection methods are often compared at small feature limits, where a score mixes two questions: whether a method returns any subset within the limit, and how well its subsets predict. We separate the two for three NSGA-II configurations (an opposition-based start, a sparse start, and our implementation of NSGAII-MIIP) on 46 datasets with 16 to 192{,}000 features and up to three classifiers, plus a four-dataset pilot. With a limit of ten features, reach accounts for 98 to 100 percent of the area-under-the-front difference between the two starts. At sizes both starts reach, mean test F1 differences stay within 0.02 under two classifiers, and under the third only without the smallest test sets. At benchmark budgets, the opposition-based start never reaches ten features beyond 512 features; with ten times the budget, it does on all four pilot datasets. Such comparisons should report reach, scores at matched sizes, and the budget.

\vspace{6pt}\textbf{Keywords:} evolutionary feature selection, multi-objective optimization, initialization, cardinality coverage, evaluation budget.
\end{abstract}
